\section{Exactness and unipotent completion}\label{sec:mhs}

We prove Theorem~\ref{thm:completion}. Ordinary fundamental-group
exactness identifies the fibre image with the commutator at each finite
unipotent stage; the mixed Hodge structures then force this image to be
central.
The proof keeps these finite stages canonical so that every map to
which a weight argument is applied is a mixed-Hodge morphism.

\subsection{The canonical quotients and their weights}
Fix a connected smooth complex quasi-projective variety $V$ and a base
point $v\in V$. Put $\Gamma=\pi_1(V,v)$ and write $\m_V$ for the
complete pronilpotent Lie algebra of its rational unipotent completion.
In characteristic zero the exponential and logarithm identify a
unipotent algebraic group with its nilpotent Lie algebra, using the
Baker--Campbell--Hausdorff multiplication. For $c\geq1$, define the
\emph{canonical $c$-step quotient} by
\begin{equation}\label{eq:canonical-quotient}
 \g_c=\m_V/\cl{\gamma_{c+1}\m_V},
 \qquad U_c=\exp(\g_c),
\end{equation}
where $\gamma_1\m_V=\m_V$,
$\gamma_{r+1}\m_V=[\m_V,\gamma_r\m_V]$, and closure is taken in the
inverse-limit topology. The group $\Gamma$ is finitely generated, so
each $\g_c$ is finite-dimensional. The canonical image of $\Gamma$ is
Zariski dense in $U_c$, and
\begin{equation}\label{eq:hurewicz}
 \g_c/[\g_c,\g_c]\simeq H_1(V,\Q).
\end{equation}
These constructions and their universal properties can be obtained from
the completed group algebra
\cite[\S2.5, Theorem 2.5.3 and Proposition 2.5.5]{Hain1987}.

We recall exactly the mixed-Hodge information used below. A rational
mixed Hodge structure is a finite-dimensional rational vector space
with an increasing weight filtration $W$ and a decreasing Hodge
filtration on its complexification, such that each weight-graded piece
is a pure Hodge structure. We say it has \emph{weights at most $r$} if
$W_r$ is the entire space, and is \emph{pure of weight $r$} if also
$W_{r-1}=0$. Morphisms preserve the filtrations and are strict: their
images receive the same filtration as subobjects of the target and as
quotients of the source. The category is abelian
\cite[Theorem 2.3.5 and \S2.3.8]{Deligne1971}. Tensor weights add.

For smooth quasi-projective $V$, first cohomology has weights $1,2$,
so $H_1(V,\Q)$ has weights $-1,-2$
\cite[Corollary 3.2.15(ii)]{Deligne1971}. The canonical Lie quotients
$\g_c$ carry mixed Hodge structures for which brackets and the
Hurewicz map \eqref{eq:hurewicz} are morphisms. Moreover a pointed
algebraic map induces a morphism of these structures. We use Hain's
pointed construction \cite[Theorem 6.3.1, p.~321]{Hain1987} for this
functoriality. It equips the completed Lie algebra with a pro-mixed
Hodge structure compatible with its bracket. The closed lower-central
ideals are topologically generated by iterated bracket images and are
therefore pro-mixed-Hodge subobjects. The finite-dimensional quotients
\eqref{eq:canonical-quotient} inherit ordinary mixed Hodge structures;
their transition maps and the maps induced by pointed algebraic
morphisms are morphisms of these structures. Morgan's foundational construction
\cite{Morgan1978} is part of the history of the method, but its original
functoriality assertion was corrected in \cite{Morgan1986}; the stronger
pointed statement needed here is supplied by Hain.

At finite stages, all lower-central ideals, images, kernels, and
abelianizations below are therefore ordinary subobjects or quotients in
the abelian category of mixed Hodge structures. Arbitrary unipotent
representation quotients need not themselves carry the required Hodge
structure. We prove the bound on the canonical quotients first and only
then pass to arbitrary representations by universality.

We now assume all the hypotheses of Theorem~\ref{thm:completion}, with
$a:V\to A$ dominant, $F$ a smooth connected general fibre, and the
base point $v$ chosen in $F$. We use the ordinary exact sequence
\eqref{eq:exact}; properness and specialness of the fibre play no role.

\subsection{The fibre image is the commutator at every unipotent stage}
Let
\[
 N=\im\bigl(\pi_1(F,v)\longrightarrow\Gamma\bigr).
\]
By \eqref{eq:exact}, $N$ is normal and equals $\ker a_*$. The fundamental
group of a semiabelian variety is free abelian. The quasi-Albanese map
induces an isomorphism
\begin{equation}\label{eq:alb-integral-h1}
 H_1(V,\Z)/\Tor H_1(V,\Z)\xrightarrow{\sim}H_1(A,\Z),
\end{equation}
as in \cite[Lemmas 3.11--3.12, arXiv v2 PDF, pp.~20--21]{Fujino2024}.
The first of those lemmas states integral surjectivity with torsion
kernel; the proof of the second gives the rational mixed-Hodge identification. Thus
\begin{equation}\label{eq:almost-derived}
 [\Gamma,\Gamma]\subseteq N,
 \qquad
 N/[\Gamma,\Gamma]\simeq\Tor H_1(V,\Z).
\end{equation}
In particular the latter quotient is finite. We do not assert equality
of the two discrete subgroups when this torsion is nonzero.

Fix $c$. Let $K_c$ be the Zariski closure of the image of $N$ in $U_c$,
and let $\kk_c=\Lie K_c$. The pointed fibre inclusion induces a map of
canonical $c$-step Lie quotients
\[
 \g_c(F)\longrightarrow\g_c(V).
\]
Its image is $\kk_c$: images of morphisms of unipotent algebraic groups
are closed, and the fibre group is dense in its own completion. This
also equips $\kk_c$ with its image mixed Hodge structure.

We claim
\begin{equation}\label{eq:k-derived}
 K_c=[U_c,U_c],\qquad\kk_c=[\g_c,\g_c].
\end{equation}
Indeed, every element of $N$ has a positive power in $[\Gamma,\Gamma]$
by \eqref{eq:almost-derived}. The abelianization $U_c/[U_c,U_c]$ is an
additive vector group, which has no nonidentity torsion. Hence the image
of $N$ in it is zero and $K_c\subseteq[U_c,U_c]$. Conversely, normality
of $N$ and density of $\Gamma$ imply normality of $K_c$ in $U_c$. The
dense image of $\Gamma$ in $U_c/K_c$ is abelian, because
$[\Gamma,\Gamma]\subseteq N$. The commutator morphism of the algebraic
quotient vanishes on this dense subset of its product, and so the
quotient is abelian. This gives the opposite inclusion and proves
\eqref{eq:k-derived}. In particular, no left-exactness assertion about
unipotent completion has been used.

\subsection{Two weight bounds force centrality}
We continue with the fixed finite-dimensional nilpotent Lie algebra
$\g_c$. Its lower-central graded pieces are generated by iterated
brackets of its abelianization. More precisely, the $r$-fold iterated
bracket induces a surjective mixed-Hodge morphism
\[
 H_1(V,\Q)^{\otimes r}\longrightarrow
 \gamma_r\g_c/\gamma_{r+1}\g_c.
\]
The target therefore has weights at most $-r$. Starting with
$\gamma_{c+1}\g_c=0$ and taking successive extensions down the
lower-central filtration shows that $\gamma_r\g_c$ itself has weights
at most $-r$. In particular,
\begin{equation}\label{eq:negative-weights}
 W_{-1}\g_c=\g_c,
 \qquad \kk_c=[\g_c,\g_c]\subseteq W_{-2}\g_c.
\end{equation}
Consider its fibre-image abelianization
\[
 B_c=\kk_c/[\kk_c,\kk_c].
\]
The surjection $\g_c(F)\to\kk_c$ induces a surjective mixed-Hodge map
\begin{equation}\label{eq:fibre-h1-onto}
 H_1(F,\Q)\longrightarrow B_c.
\end{equation}
The group $H_1(F,\Q)$ has only weights $-1,-2$, because the chosen fibre
is smooth and quasi-projective. Strictness of
\eqref{eq:fibre-h1-onto} gives $W_{-3}B_c=0$. On the other hand,
strictness for the image $\kk_c\subseteq\g_c$ and
\eqref{eq:negative-weights} give $W_{-2}B_c=B_c$. Therefore
\begin{equation}\label{eq:pure-minus-two}
 B_c\text{ is pure of weight }-2.
\end{equation}
This is the point where geometry of the fibre enters: its first
homology forbids weights below $-2$, even though it maps into a
commutator ideal.

Because $\kk_c$ is an ideal, the bracket followed by its abelianization
is a mixed-Hodge morphism
\[
 \g_c\otimes\kk_c\longrightarrow B_c.
\]
The source has weights at most $-3$, and the target has $W_{-3}=0$ by
\eqref{eq:pure-minus-two}. The morphism must vanish. Thus
\begin{equation}\label{eq:action-zero}
 [\g_c,\kk_c]\subseteq[\kk_c,\kk_c].
\end{equation}
Substituting \eqref{eq:k-derived} and using the Jacobi lower-central
estimate yields
\[
 \gamma_3\g_c
 \subseteq[\gamma_2\g_c,\gamma_2\g_c]
 \subseteq\gamma_4\g_c
 \subseteq\gamma_3\g_c.
\]
Hence $\gamma_3\g_c=\gamma_4\g_c$. Repeated bracketing with $\g_c$
makes all subsequent lower-central terms equal. Since $\g_c$ is
nilpotent, they must be zero. We have proved
\begin{equation}\label{eq:all-c-zero}
 \gamma_3\g_c=0\qquad\hbox{for every }c\geq1.
\end{equation}

The complete Lie algebra $\m_V$ is the separated inverse limit of these
canonical quotients. Every element of $\gamma_3\m_V$ has zero image in
every $\g_c$ by \eqref{eq:all-c-zero}; separatedness therefore gives
$\gamma_3\m_V=0$. Its closure is zero as well. By the definition
\eqref{eq:canonical-quotient}, the map $\m_V\to\g_2$ is consequently
an isomorphism. Thus the canonical tower actually stabilizes at $c=2$,
and the completion is finite-dimensional. Its abelianization is
$H_1(V,\Q)$ and its commutator is a quotient of
$\bigwedge^2H_1(V,\Q)$. The Baker--Campbell--Hausdorff correspondence
gives the same class bound for the unipotent algebraic group.

It remains to pass to complex unipotent representations. For the
finitely generated group $\Gamma$, unipotent completion commutes with
extension of the characteristic-zero ground field \cite[\S4]{Hain2015}. More explicitly, the completion over
$\C$ is the scalar extension of the rational completion and has the
universal property for homomorphisms into complex unipotent groups.
After extension from $\Q$ to $\C$, this universal property
therefore bounds every complex unipotent representation image. This
completes the proof of Theorem~\ref{thm:completion}.

\begin{remark}\label{rem:proper-fibre}
If $H_1(F,\Q)$ is pure of weight $-1$, in particular if $F$ is proper,
the same proof gives $B_c=0$. A nilpotent Lie algebra with zero
abelianization is zero, so $\kk_c=0$ and the completion is abelian.
This reflects the difference between complete fibres and fibres whose
first homology has a boundary contribution of weight $-2$. For smooth
quasi-projective varieties whose quasi-Albanese map is proper and
surjective, Aguilar Aguilar--Campana proved that every torsion-free
nilpotent quotient is abelian
\cite[arXiv version, Corollary 2(2)]{AguilarAguilarCampana2025}.
\end{remark}

\begin{remark}
The zero cases cause no exception. If $A$ is a point, then
$H_1(V,\Q)=0$ and all the canonical nilpotent Lie quotients vanish. If
$F$ is a point, exactness identifies $\Gamma$ with the free abelian
group $\Lambda$. Neither conclusion requires division by a positive
Betti number.
\end{remark}
